% Bewertungspaket zum Gesamttest «Quadratische Funktionen» — Quelle des PDFs (python3 scripts/build-lp-pdf.py).
% Ganze Punkte je Teilschritt; Werte mit python3 nachgerechnet (02.10.2026).
\documentclass[11pt]{article}
\usepackage{../lp-druck}
\renewcommand{\lpfuss}{Bewertungspaket}
\newcolumntype{P}{>{\centering\arraybackslash}p{6mm}}
\newenvironment{raster}{\par\smallskip\noindent\tabularx{\linewidth}{@{}>{\raggedright\arraybackslash}p{38mm}P X@{}}\toprule
  \small\textsc{Teilschritt} & \small\textsc{P} & \small\textsc{Musterlösung}\\\midrule}{\bottomrule\endtabularx\par}
\newcommand{\fehler}[1]{\par\smallskip{\small\textcolor{lprot}{\bfseries Typische Fehler:} #1}}
\begin{document}
\lpkopf{Bewertungspaket zum Gesamttest}{}

\begin{lpachtung}{Erst lösen, dann öffnen}
Dieses Paket enthält die Musterlösung. Wer es vor dem Lösen liest, misst am Ende nur, wie gut das Abschreiben gelingt.
\end{lpachtung}

\section*{1 · So gehst du vor}
\begin{enumerate}[leftmargin=*,itemsep=2pt]
\item \textbf{Gesamttest lösen} — vollständig, mit Rechenweg, ohne dieses Paket.
\item \textbf{Fotografieren oder scannen}, gut lesbar, eine Seite pro Bild. \textbf{Kein Name} auf den Bildern.
\item \textbf{Bewerten lassen.} Ob du dafür eine KI nutzt und welche, entscheidest du selbst und in eigener Verantwortung —
  je nachdem, was dir aufgrund deines Alters und deiner persönlichen Situation erlaubt ist (Altersgrenzen und
  Nutzungsbedingungen des Anbieters, allenfalls Einverständnis deiner Eltern). Lade nur deine Lösung und dieses PDF hoch,
  keine persönlichen Daten, und füge den Auftrag aus Abschnitt~2 ein. Ohne KI kannst du dich mit Abschnitt~3 selbst bewerten.
\item \textbf{Rückmeldung prüfen.} Stimmt, was die KI aus deiner Handschrift gelesen hat? Eine KI kann sich irren —
  im Zweifel gilt das Raster in Abschnitt~3.
\item \textbf{Gezielt wiederholen} nach Abschnitt~4.
\end{enumerate}

\section*{2 · Auftrag an die KI}
\begin{tcolorbox}[colback=lplinie!25,colframe=lplinie,boxrule=0.5pt,arc=1mm,fontupper=\small]
Du bist Korrektorin oder Korrektor für Mathematik an einer Schweizer Berufsmaturitätsschule. Im Anhang findest du (1)~das Bewertungspaket zum Gesamttest «Quadratische Funktionen» mit Musterlösung und Punkteraster und (2)~Fotos meiner handschriftlichen Lösung.\par\medskip
Geh so vor:\par
1. Schreib zu jeder Aufgabe G1 bis G8 zuerst ab, was du in meiner Lösung liest — Rechenweg und Ergebnis. Was du nicht sicher lesen kannst, markierst du mit [unsicher]. Nicht raten.\par
2. Vergib die Punkte genau nach dem Raster im Bewertungspaket (Abschnitt 3). Ziehe einen Fehler nur einmal ab: Rechne ich mit einem falschen Zwischenergebnis richtig weiter, gibt es die Folgepunkte.\par
3. Andere richtige Lösungswege zählen voll. Dezimalpunkt und Dezimalkomma sind beide richtig. Ein Ergebnis ohne nachvollziehbaren Weg gibt höchstens den Ergebnispunkt.\par
4. Begründe jeden Abzug in einem Satz und nenne die Fehlerart (zum Beispiel «Vorzeichen in der Klammer»). Schreib die Musterlösung nicht ab — zeig mir, wo mein Fehler liegt.\par
5. Gib zum Schluss eine Tabelle «Aufgabe | Punkte | Begründung», die Summe von 24 und — nach der Selbsteinschätzung in Abschnitt 4 — welches Kapitel des Leitprogramms ich wiederholen soll. Keine Note.\par
6. Fehlt ein Foto oder ist eine Aufgabe unlesbar, sag es, statt sie zu bewerten.
\end{tcolorbox}

\section*{3 · Musterlösung und Punkteraster}
\textbf{Allgemein:} Ein Fehler wird nur einmal abgezogen (Folgefehler). Gleichwertige Wege zählen voll.
Ergebnis ohne Weg: höchstens der Ergebnispunkt. Teil A und C sind ohne Taschenrechner gestellt.

\begin{lpaufgabe}{G1}{$f(x) = -(x-2)^2 + 9$}{4 P}
\begin{raster}
Scheitelpunkt & 1 & $S(2 \mid 9)$\\
Öffnung & 1 & nach unten, da $a = -1 < 0$\\
Nullstellen mit Weg & 1 & $(x-2)^2 = 9 \Rightarrow x - 2 = \pm 3 \Rightarrow x_1 = -1,\ x_2 = 5$ (beide nötig)\\
Produktform & 1 & $f(x) = -(x+1)(x-5)$\\
\end{raster}
\fehler{$S(-2 \mid 9)$ (Vorzeichen in der Klammer) $-1$, Nullstellen dazu folgerichtig voll; nur eine Nullstelle ($\pm$ vergessen) $-1$; Produktform ohne Minus vor der Klammer $-1$.}
\end{lpaufgabe}

\begin{lpaufgabe}{G2}{$f(x) = 2x^2 - 12x + 10$}{3 P}
\begin{raster}
Ausklammern bzw.\ Ansatz & 1 & $2(x^2 - 6x + 5)$ — oder $x_s = -\frac{b}{2a} = 3$, $y_s = f(3)$\\
Scheitelform & 1 & $2\big((x-3)^2 - 4\big) = 2(x-3)^2 - 8$\\
Scheitel und $y$-Achse & 1 & $S(3 \mid -8)$, Schnittpunkt $(0 \mid 10)$\\
\end{raster}
\fehler{$2(x-3)^2 - 4$ (die $-4$ nicht mit 2 multipliziert) $-1$, Scheitel dazu folgerichtig; $(0 \mid 10)$ fehlt $-1$.}
\end{lpaufgabe}

\begin{lpaufgabe}{G3}{Graph $\to$ Scheitelform}{2 P}
\begin{raster}
Scheitel abgelesen & 1 & $S(-1 \mid -2)$\\
Streckfaktor und Gleichung & 1 & ein Schritt nach rechts $\to$ 1 hinauf, also $a = 1$: $f(x) = (x+1)^2 - 2$\\
\end{raster}
\fehler{$(x-1)^2 - 2$ (Vorzeichen in der Klammer) $-1$.}
\end{lpaufgabe}

\begin{lpaufgabe}{G4}{$f(x) = 3(x+1)(x-2)$}{2 P}
\begin{raster}
Klammern ausmultipliziert & 1 & $(x+1)(x-2) = x^2 - x - 2$\\
Faktor auf alle Glieder & 1 & $f(x) = 3x^2 - 3x - 6$\\
\end{raster}
\fehler{$3x^2 - x - 2$ (Faktor nur aufs erste Glied) $-1$.}
\end{lpaufgabe}

\begin{lpaufgabe}{G5}{Scheitel $S(3 \mid -2)$, Punkt $P(5 \mid 6)$}{3 P}
\begin{raster}
Ansatz Scheitelform & 1 & $f(x) = a(x-3)^2 - 2$\\
$a$ bestimmt & 1 & $6 = a \cdot 4 - 2 \Rightarrow a = 2$\\
Gleichung & 1 & $f(x) = 2(x-3)^2 - 2$ (gleichwertig $2x^2 - 12x + 16$)\\
\end{raster}
\fehler{Ansatz $a(x+3)^2 - 2$ $-1$, Rest folgerichtig voll.}
\end{lpaufgabe}

\begin{lpaufgabe}{G6}{Wasserstrahl}{3 P}
\begin{raster}
Ansatz & 1 & Nullstellen 0 und 8: $h(x) = a\,x(x-8)$ — oder Scheitel $(4 \mid 4)$: $a(x-4)^2 + 4$\\
$a$ bestimmt & 1 & $4 = a \cdot 4 \cdot (-4) \Rightarrow a = -0.25$\\
Gleichung & 1 & $h(x) = -0.25\,x(x-8) = -0.25x^2 + 2x$\\
\end{raster}
\fehler{$a > 0$ (nach oben geöffnet, Widerspruch zum Sachverhalt) $-1$; Scheitel bei $(8 \mid 4)$ statt in der Mitte $-1$.}
\end{lpaufgabe}

\begin{lpaufgabe}{G7}{$h(t) = -5t^2 + 20t + 1.5$}{3 P}
\begin{raster}
Zeitpunkt mit Weg & 1 & $t = -\frac{b}{2a} = -\frac{20}{2 \cdot (-5)} = 2$ — oder quadratisch ergänzen\\
Höhe & 1 & $h(2) = -20 + 40 + 1.5 = 21.5$\\
Antwort mit Einheiten & 1 & Nach 2\,s ist der Ball am höchsten, 21.5\,m hoch.\\
\end{raster}
\fehler{$t = -2$ (Minus in $-\frac{b}{2a}$) $-1$, Höhe dazu folgerichtig; nur «2» ohne Höhe oder ohne Einheiten $-1$.}
\end{lpaufgabe}

\begin{lpaufgabe}{G8}{Regenrinne aus 24\,cm Blech}{4 P}
\begin{raster}
Zielfunktion & 1 & Boden $24 - 2x$: $A(x) = x(24 - 2x)$\\
Stelle des Maximums & 1 & Nullstellen 0 und 12, Mitte $x = 6$ — oder $-\frac{b}{2a}$ aus $-2x^2 + 24x$\\
Grösste Fläche & 1 & $A(6) = 6 \cdot 12 = 72$\\
Antwort & 1 & Seiten 6\,cm hoch, Boden 12\,cm, Querschnitt 72\,cm$^2$\\
\end{raster}
\fehler{$A(x) = x(24 - x)$ (nur eine Seite abgezogen) $-1$, Rest folgerichtig (dann $x = 12$, $A = 144$); nur $x = 6$ ohne Fläche $-1$.}
\end{lpaufgabe}

\section*{4 · Selbsteinschätzung}
\begin{tabularx}{\linewidth}{@{}l X@{}}\toprule
21 – 24 P & Sitzt. Die Kompetenzen K1–K4 von Teilgebiet 3.3 hast du im Griff.\\
16 – 20 P & Den schwächsten Teil nochmals: Simulation und Übungen des Kapitels.\\
10 – 15 P & Teil A (G1–G4) $\to$ Kapitel 1–3, Teil B (G5–G6) $\to$ Kapitel 4, Teil C (G7–G8) $\to$ Kapitel 5.\\
0 – 9 P & Zurück zu Kapitel 1.\\\bottomrule
\end{tabularx}
\end{document}
